Language models in the low-billions parameter range are increasingly deployed on edge devices, where fixed compute budgets must continuously absorb new domain tasks via parameter-efficient adapters (e.g., LoRA) without catastrophic forgetting. Continual learning typically restricts updates to a parameter subset chosen by gradient attribution---using a calibration backward pass to score importance. We show that previous evaluations of this approach suffer from a fundamental confound: under standard zero-initialized LoRA ($B=0$), the gradient with respect to $A$ is identically zero ($\partial\mathcal{L}/\partial A=0$), forcing gradient masks to select strictly from the $B$ matrix, whereas standard random baselines draw from both $A$ and $B$. 
The comparison these evaluations run is itself null 
% (Backward Transfer $0.437 \pm 0.061$ versus $0.428 \pm 0.019$, $t(4)=0.41$, $p=0.70$)
; matching the parameter pool is what surfaces the effect. 
Isolating the selection criterion by restricting both approaches to the same parameter pool 
% ($10\%$ of LoRA $B$, 1.65M parameters)
, we find a negative result: a uniform random mask consistently achieves lower loss-based backward transfer than gradient attribution across seven seeds on Qwen2.5-3B 
% (mean BWT $0.294$ vs.\ $0.449$, $\Delta = -0.155$ [Random $-$ Grad] on the V100 stack, $p = 0.0011$; 7/7 seeds)
. The effect is concentrated on the medical task, which contributes 80\% of the aggregate; code and math move in the same direction, and legal shows no reliable difference. This advantage is consistent with a stability--plasticity trade-off, in which random masking slightly reduces single-task fitting. On accuracy the two masks are statistically indistinguishable on the two discrete-option tasks (Medical, Legal) 
% ($+0.94$\,pp, 95\% CI $[-2.91, +4.80]$)
. 
% A 4-percentage-point non-inferiority margin specified in advance is not met (upper bound $+4.80$\,pp). 
% On those same two tasks, accuracy differences of up to about 4.8 percentage points in the gradient mask's favour are not excluded by this measurement.
On loss-based forgetting, which is the metric this paper measures, gradient attribution is measurably worse than density-matched random masking, and the attribution pass---adding per-task backward overhead and calibration storage---can be dropped in favour of a simpler, zero-cost alternative for resource-constrained LLM continual learning.
